% Presentation of original arrow notation. See flow-manifest.json.
\begin{center}
\begin{tikzpicture}[book flow node/.style={draw=BookHeading,line width=.65pt,fill=BookPanel,rounded corners=1.5pt,align=center,inner sep=5pt,minimum height=10mm,font=\fontsize{9}{11}\selectfont,execute at begin node={\hyphenpenalty=10000\exhyphenpenalty=10000}},book flow edge/.style={-{Latex[length=2mm]},draw=BookHeading,line width=.7pt}]
\node[book flow node,text width=48mm] (n0) at (-3,0) {semantic.\allowbreak{}accent \#7aa7e8 *};
\node[book flow node,text width=48mm] (n1) at (3,0) {border.\allowbreak{}secondary};
\node[book flow node,text width=48mm] (n2) at (3,-1.6) {ui.\allowbreak{}cursor};
\node[book flow node,text width=48mm] (n3) at (3,-3.2) {ui.\allowbreak{}focus\allowbreak{}Gutter};
\node[book flow node,text width=48mm] (n4) at (-3,-5) {surface.\allowbreak{}primary.\allowbreak{}bg \#171827};
\node[book flow node,text width=48mm] (n5) at (3,-5) {surface.\allowbreak{}primary};
\node[book flow node,text width=48mm] (n6) at (3,-6.6) {ui.\allowbreak{}focus\allowbreak{}Gutter.\allowbreak{}bg};
\draw[book flow edge] (n0) -- (n1);
\draw[book flow edge] (n0) -- (n2);
\draw[book flow edge] (n0) -- (n3);
\draw[book flow edge] (n4) -- (n5);
\draw[book flow edge] (n4) -- (n6);
\end{tikzpicture}
\end{center}
